\documentclass[11pt]{article}

\usepackage[T1]{fontenc}
\usepackage{lmodern}
\usepackage[margin=1in]{geometry}
\usepackage{amsmath,amssymb,amscd,amsthm}
\usepackage{booktabs,tabularx}
\usepackage{microtype}
\usepackage[backend=biber,style=numeric,sorting=none]{biblatex}
\usepackage[hidelinks]{hyperref}
\hypersetup{pdftitle={Compatibility and Affine Witness Disclosure}}
\addbibresource{references.bib}

\newtheorem{proposition}{Proposition}[section]
\DeclareMathOperator{\rank}{rank}
\DeclareMathOperator{\TAUT}{TAUT}
\newcommand{\Ftwo}{\mathbb F_2}
\newcommand{\NP}{\mathsf{NP}}
\newcommand{\PT}{\mathsf{P}}
\newcommand{\coNP}{\mathsf{coNP}}

\begin{document}
\begin{center}
  {\LARGE Compatibility and Affine Witness Disclosure\par}
  \smallskip
  \textit{A study of collision squares and their realizations}
\end{center}

\begin{abstract}
A comparison map defines a collision relation on its source. An operation
descends through a surjective comparison exactly when it preserves that
relation. We distinguish exact coding identities, provable equivalence,
semantic quotients, and proof translations with resource bounds. For the
fresh-sign Harmonic Cochain Problem, we identify the collision space after
an affine disclosure with an affine witness fiber times a kernel. This
comparison preserves the diagonal projector and gives an exact count of
the surviving witnesses.
\end{abstract}

\section{Compatibility, descent, and faithful realization}
\label{sec:compatibility}

For a map \(q\colon A\to Q\), a collision is a pair \((a,a')\) with
\(q(a)=q(a')\). The collision space \(A\times_Q A\) is the pullback of the
diagonal of \(Q\) along \((q,q)\). The diagonal of \(A\) is its distinguished
subset, and the residual consists of the distinct pairs with equal images.
Every comparison below specifies the equality being imposed and the objects
on which it is imposed.

\begin{proposition}[Descent along a comparison]\label{prop:descent}
Let \(q\colon A\twoheadrightarrow Q\) be surjective and \(f\colon A\to Z\).
There is a unique \(\bar f\colon Q\to Z\) with \(\bar f q=f\) if and only if
\[
  q(a)=q(a')\quad\Longrightarrow\quad f(a)=f(a').
\]
In particular, \(D\colon A\to A\) descends to an endomorphism of \(Q\)
exactly when \((D,D)\) preserves \(A\times_Q A\).
\end{proposition}

\begin{proof}
If \(\bar f q=f\), apply \(\bar f\) to the equality of images.
Conversely, define \(\bar f(x)=f(a)\) using a representative \(a\) with
\(q(a)=x\). The collision condition makes the value independent of the
representative. Surjectivity gives uniqueness. For the endomorphism
statement, take \(f=qD\).
\end{proof}

A realization \(E\colon V\to Z\) with a left inverse \(J\) reflects
equality. Applying \(J\) to \(E(v)=E(v')\) gives \(v=v'\). If \(V\) is
itself a quotient, this detects equality of its classes. The descent
condition concerns an operation on presentations, while the left inverse
concerns the information retained by a realization.

\subsection{The coding and provability squares}

Fix a G\"odel coding \(g\) with tracked numeral substitution. Let
\(\mathrm{Form}_1\) consist of formulas with at most one free variable,
designated \(x\), and let \(\mathrm{Sent}\) be the sentences. The restrictions
\(g_1\) and \(g_0\) take values in their respective images
\(\mathrm{Code}_1\) and \(\mathrm{Code}_0\). With \(\overline n\) the numeral
for \(n\), define
\[
  \delta_g(\varphi)=\varphi[\overline{g_1(\varphi)}/x],
  \qquad d_g(n)=\operatorname{sub}_g(n,n).
\]
The numeral substitution law gives the exact identity
\(d_g g_1=g_0\delta_g\) \cite{grabmayr,goldfarb}.
\begin{equation}\label{eq:coding}
\begin{CD}
  \mathrm{Form}_1 @>{\delta_g}>> \mathrm{Sent}\\
  @V{g_1}VV @VV{g_0}V\\
  \mathrm{Code}_1 @>{d_g}>> \mathrm{Code}_0 .
\end{CD}
\end{equation}

For a theory \(T\) with sufficient arithmetic and a unary formula
\(\theta(z)\), the diagonal lemma supplies a sentence \(G\) with
\[
  T\vdash G\leftrightarrow\theta(\overline{g_0(G)})
  \qquad\text{\cite{goldfarb,openlogic}}.
\]
Writing \(q_T\colon\mathrm{Sent}\to H_T\) for the quotient by provable
equivalence and \(h_\theta(n)=[\theta(\overline n)]_T\), this is the equality
\(q_T(G)=h_\theta(g_0(G))\) in \(H_T\). It is an agreement at the
constructed sentence \(G\). The coding square \eqref{eq:coding} instead
holds for every formula in its domain. To descend own-code substitution
through an equivalence on formulas, one must also show that equivalent
formula presentations yield equivalent output sentences, as in
Proposition~\ref{prop:descent}.

\subsection{Provability and selected semantics}

For a family \(\mathcal M\) of models of \(T\), let \(Q_{\mathcal M}\)
be the quotient of sentences by equivalence in every model in the family.
Soundness gives a surjective Boolean ring homomorphism
\(c_{\mathcal M}\colon H_T\to Q_{\mathcal M}\), with
\[
 Q_{\mathcal M}\cong H_T/\ker c_{\mathcal M}.
\]
Writing \(I_{\mathcal M}=\ker c_{\mathcal M}\), the map
\((x,y)\mapsto (x,y-x)\) gives a bijection
\[
 H_T\times_{Q_{\mathcal M}}H_T\;\cong\; H_T\times I_{\mathcal M}.
\]
Its inverse adds the difference to the first class, and equal pairs
correspond to the zero element of the ideal.

For a sound, computably axiomatized extension of Robinson arithmetic,
take \(\mathcal M=\{\mathbb N\}\). A true unprovable G\"odel sentence
gives distinct classes \([G]_T\) and \([\top]_T\) with equal semantic image.
In the corresponding ideal, \([\neg G]_T\) is nonzero. When
\(\mathcal M\) contains all models of \(T\), first-order completeness
makes the comparison injective and the residual empty
\cite[Chapters 23 and 37]{openlogic}.

\section{Affine witness disclosure}
\label{sec:disclosure}

In Paper~1's fresh-sign family, fix an integer matrix \(A\) with
\(A^{\mathsf T}A=dI\), \(d>0\), a public integer \(a\), and a sign mask
\(\tau\). The observed cochain and the orthogonal residual for message
signs \(s\) are
\[
 b=(a,0,(d+1)A\tau),\qquad
 r_s=(a,-ds,A\tau)
 \qquad\text{\cite{paperone}}.
\]
For key and message bits \(u,v\), with sign \(2v_j-1\) for \(v_j\),
compatibility is \(u+v=m\) in \(\Ftwo^n\), where \(m_j=1\) for a
negative mask sign. Adding explicitly given affine constraints
\(H_u u+H_v v=c_0\) and substituting \(u=m+v\) gives the witness space
\[
 Q=H_u+H_v,\qquad c=c_0+H_um,\qquad
 W=\{v\in\Ftwo^n:Qv=c\},\qquad K=\ker Q .
\]
Every \(v\in W\) determines one compatible key. Disclosing information
coordinates \(J\) of the residual reveals the corresponding message bits
since \(r_{s,j}=d(1-2v_j)\) in that block. The disclosed bits define the
projection \(D\colon\Ftwo^n\to\Ftwo^J\).

For key signs \(\kappa\), the corresponding closest-vector lattice is
generated by \(G_\kappa z=(0,z,A(\kappa\odot z))\), where \(\odot\) is
entrywise multiplication. Since \(\kappa\odot s=\tau\),
\(b-G_\kappa(ds)=r_s\). The identity \(A^{\mathsf T}A=dI\) gives
\(G_\kappa^{\mathsf T}G_\kappa=(d+1)I\) and
\(G_\kappa^{\mathsf T}r_s=0\). Thus for \(z\in\mathbb Z^n\),
\begin{equation}\label{eq:cvp-cost}
  \|b-G_\kappa z\|^2
    =a^2+d(d+1)n+(d+1)\|z-ds\|^2.
\end{equation}
Each compatible key gives a unique closest vector in its lattice. The
minimum is common across keys and the residual vector retains the
message in its information block.

\begin{proposition}[Disclosure collision square]\label{prop:disclosure}
Put \(K_J=K\cap\ker D\). The maps \(\Psi\) and \(\Psi_J\) are bijections
and the square is a pullback. Both take \((v,v')\) to \((v,v'-v)\), marking
the diagonal by zero direction.
\begin{equation}\label{eq:disclosure}
\begin{CD}
  W\times_D W @>{\Psi_J}>> W\times K_J\\
  @VVV @VVV\\
  W\times W @>{\Psi}>> W\times K .
\end{CD}
\end{equation}
\end{proposition}

\begin{proof}
The difference of two solutions of \(Qv=c\) lies in \(K\). Adding a
direction in \(K\) to a solution gives a solution, which gives the inverse
of \(\Psi\). Since \(D(v'-v)=D(v')-D(v)\), the equality of disclosures
cuts out exactly the preimage of \(W\times K_J\). Equal witnesses have
zero difference, and a zero difference forces equality.
\end{proof}

\subsection{Boolean coordinate rings}
\label{sec:boolean-disclosure}

For a finite set \(T\), write \(\mathcal A(T)=\Ftwo^T\) for its function
algebra. Evaluation identifies
\(\Ftwo[t_1,\ldots,t_n]/(t_i^2-t_i)_i\) with
\(\mathcal A(\Ftwo^n)\). Indeed, the polynomials
\(\prod_{b_i=1}t_i\prod_{b_i=0}(1-t_i)\), indexed by \(b\in\Ftwo^n\),
give all point indicators. Reduction to square-free monomials bounds
the dimension by \(2^n\), so evaluation is an isomorphism. In this
product of fields, quotienting by a family of equations keeps exactly
the factors at their common zeros.

The collision space \(C_J=W\times_D W\) therefore has coordinate ring
\begin{align}
\mathcal A(C_J)
 &\cong
 \frac{\Ftwo[v,v']}{(v_i^2-v_i,\ (v_i')^2-v_i',\ Qv-c,\ Qv'-c,\
                     D(v'-v))} \notag\\
 &\cong
 \frac{\Ftwo[v,w]}{(v_i^2-v_i,\ w_i^2-w_i,\ Qv-c,\ Qw,\ Dw)}
 \cong \mathcal A(W\times K_J).
 \label{eq:collision-ring}
\end{align}
The middle isomorphism uses \(w=v'-v\), with inverse \(v'=v+w\).
Subtracting the two affine equations gives \(Qw=0\), and the Boolean
equations are preserved because squaring is additive over \(\Ftwo\).
Under this change of coordinates the diagonal indicator becomes
\(e_\Delta(w)=\prod_i(1-w_i)\). Its complementary idempotent
\(1-e_\Delta\) marks the distinct pairs with equal disclosures.
The two projectors split the coordinate ring as
\[
  \mathcal A(C_J)\cong
  \mathcal A(W)\times\mathcal A\bigl(W\times(K_J\setminus\{0\})\bigr).
\]
On taking spectra the two factors give complementary open and closed
subschemes \cite[Tag 04PN]{stacksclopen}.

\paragraph{First-colon residual.}
Let \(S=\Ftwo[v,w]\), let \(I_C\) be the ideal of collision equations
in \eqref{eq:collision-ring}, and put \(I_\Delta=I_C+(w_1,\ldots,w_n)\).
In \(C=S/I_C\), the diagonal ideal \(J_\Delta=I_\Delta/I_C\) consists of
the functions vanishing on the diagonal, so \(J_\Delta=(1-e_\Delta)C\).
An element of \(S\) lies in \(I_C:I_\Delta\) exactly when its image in
\(C\) annihilates \(J_\Delta\). It follows that
\[
 S/(I_C:I_\Delta)\cong C/\operatorname{Ann}_C(J_\Delta)
 \cong \mathcal A\bigl(W\times(K_J\setminus\{0\})\bigr),
\]
since \(\operatorname{Ann}_C(J_\Delta)=e_\Delta C\). Thus the first-colon
quotient realizes the same residual as the complementary idempotent.

\subsection{Fiber size and selection}

For an attained disclosure \(y\), let \(W_y=\{v\in W:Dv=y\}\). Translation
by any \(v_0\in W_y\) identifies \(K_J\) with \(W_y\). Rank--nullity for
\(D|_K\) then gives the exact count
\[
  |W_y|=2^{\dim K-\rank(D|_K)} .
\]
The full ordered collision space has \(|W||K_J|\) pairs. Of these,
\(|W|(|K_J|-1)\) have distinct witnesses. An attained disclosure
determines the witness exactly when \(K_J=\{0\}\). Under a uniform
choice from \(W\), conditioning on an attained \(y\) is uniform on
\(W_y\), and the optimal exact recovery probability is \(1/|W_y|\).
The squared norm in \eqref{eq:cvp-cost} at \(z=ds\) descends through
every disclosure because it is constant across \(W\).

\paragraph{Two-bit example.}
Take \(A=\left(\begin{smallmatrix}1&1\\1&-1\end{smallmatrix}\right)\),
\(d=2\), \(a=3\) and \(\tau=(-1,-1)\). On the public cochain
\(b=(3,0,0,-6,0)\), the side condition \(v_1+v_2=1\) leaves messages
\((1,0)\) and \((0,1)\). Their respective residuals are
\((3,-2,2,-2,0)\) and \((3,2,-2,-2,0)\), each of squared norm \(21\).
Here \(K\) is the span of \((1,1)\). The first information coordinate
distinguishes those two messages and leaves a single witness.

For given bits \(y\), row reduction of the stacked system
\((Qv,Dv)=(c,y)\) either detects inconsistency or returns a particular
solution \(v_0\) and a basis of \(K_J\). Assigning zero to each free
variable selects a witness deterministically. The selector uses
polynomial time in the size of the explicitly given system, and
\(u=m+v_0\) recovers the corresponding key. The Python implementation
uses this elimination at the decoded cochain boundary.

For nonempty \(W\), the minimum number of coordinate disclosures needed
for exact recovery is \(k=\dim K\). Fewer than \(k\) coordinates have
rank less than \(k\) on \(K\). For the upper bound, take the \(n\times k\)
matrix whose columns are a basis of \(K\). Its column rank is \(k\), so
there are \(k\) independent rows. Disclosing those coordinates makes
\(D|_K\) injective. Row reduction finds such a set of coordinates in
polynomial time from the explicit system.

\subsection{Polynomial agreement certificates}
\label{sec:certificate-disclosure}

Paper~3's local working draft fixes distinct evaluation sites
\(a_1,\ldots,a_n\) in a field \(\mathbb F\), a received word
\(w=(w_{i0},w_{i1})_{i=1}^n\), and a Reed--Solomon pair code whose encoding
polynomials have degree below \(k\), where \(1\le k\le n\)
\cite[Section~8]{paperthreeagreement}.
A certificate has multipliers \(h_0,h_1\) of degree at most \(t\) and a
response \(R\) of degree \(D\) with \(k+t\le D<n\) and
\(R(a_i)=h_0(a_i)w_{i0}+h_1(a_i)w_{i1}\) at each site.
Dividing all three polynomials by the leading coefficient of \(R\)
makes \(R\) monic and preserves these equations. The coefficients of
\(h_j=\sum_{s=0}^{t}h_{js}X^s\) and \(R=X^D+\sum_{s=0}^{D-1}r_sX^s\) then
satisfy
\begin{equation}\label{eq:certificate-system}
 \sum_{j=0}^1 w_{ij}\sum_{s=0}^{t}a_i^s h_{js}
     -\sum_{s=0}^{D-1}a_i^s r_s=a_i^D
 \quad(1\le i\le n).
\end{equation}
For encoding polynomials \(f_0,f_1\) of a codeword, the polynomial
\(R-h_0f_0-h_1f_1\) has degree \(D\) and vanishes at every site where
both received coordinates agree with the codeword. Thus \(n-D\) is a
certified lower bound on the number of sites with a whole-pair error.

For fixed \(w,t,D\), write the system as \(Mz=b\) with
\(b=(a_i^D)_i\). Its coefficient fiber \(W_{w,t,D}\) is either empty or an
affine translate of \(\ker M\). For any linear disclosure \(L\) on the
coefficients, Proposition~\ref{prop:disclosure} applies with
\(W=W_{w,t,D}\), \(Q=M\), and \(L\) as disclosure map. The collision
square takes a pair of certificates to the first certificate and a
direction in \(\ker M\cap\ker L\), with the diagonal corresponding to
zero direction. Over \(\mathbb F_q\), each attained disclosure therefore
has \(q^{\dim(\ker M\cap\ker L)}\) certificates.

Using budget \(t=D-k\) includes every smaller budget at the same degree.
Multiplying \(h_0,h_1,R\) by \(X^{E-D}\) carries a certificate at degree
\(D\) to one at any \(D\le E<n\) with budget \(E-k\). The feasible degrees
therefore form an upper interval. Binary search with row reduction at
each degree finds its first element or detects an empty interval.
At the first feasible degree, setting free variables to zero in a
fixed coefficient order selects a normalized certificate. The search
and selection take polynomially many field operations in the size of
the explicitly given data.

When the received word varies, feasibility at a fixed \(t,D\) is the
rank condition \(\rank M=\rank[M\mid b]\), a constructible condition
given by minors on each rank stratum. Over an algebraic closure \(K\),
let \(J\subset K[w,z]\) be generated by the rows of \(Mz-b\).
Elimination gives the Zariski closure of the projection to \(w\) as
\(V(J\cap K[w])\) \cite[Section~8]{paperthreeagreement}.

\paragraph{Four-site example.}
Take the sites \(a=(1,2,3,4)\) in \(\mathbb F_5\), \(k=D=1\), and \(t=0\).
For the received columns \(w_0=a,\ w_1=2a\), the coefficient matrix is
\([a,2a,-\mathbf 1]\) and the right-hand side is \(a\). In the order
\((h_{00},h_{10},r_0)\), the five certificates are
\((1-2s,s,0)\) for \(s\in\mathbb F_5\). Fixing the first coefficient
\(h_{00}=0\) selects \((0,3,0)\). Along the line of received words
\((\lambda a,2\lambda a)\) over \(K\), each nonzero \(\lambda\) has the
certificate \((\lambda^{-1},0,0)\). Thus the zero specialization belongs
to the closure. At zero the coefficient matrix has rank one, while its
augmentation has rank two, so the fiber is empty.

On the same sites, take \(w_i=(1-a_i^{-1},0)\). Its four distinct
first coordinates give three errors against a nearest constant pair.
The first feasible degree is three, giving a certified count of one.
The Python example checks both degree searches and the equations of
the selected certificates.

\section{Resource bounds on comparisons}

\subsection{Translation of proof strings}

A Cook--Reckhow proof system for propositional tautologies is a
polynomial-time computable surjection onto \(\TAUT\) from proof strings.
For two such systems \(f_i\colon \Pi_i\to\TAUT\), a polynomial-time
translation \(\tau\colon\Pi_1\to\Pi_2\) with \(f_2\tau=f_1\) gives
\begin{equation}\label{eq:proof-translation}
\begin{CD}
  \Pi_1 @>{\tau}>> \Pi_2\\
  @V{f_1}VV @VV{f_2}V\\
  \TAUT @>{\mathrm{id}}>> \TAUT .
\end{CD}
\end{equation}
The target retains the formula proved. Cook and Reckhow call this
\emph{p-simulation} and show that it transfers a polynomial bound on
proof lengths \cite[Definitions 1.3 and 1.5, Proposition 1.6]{cookreckhow}.
Indeed, polynomial-time translation has polynomial output length.

A polynomially bounded proof system for \(\TAUT\) exists exactly when
\(\NP=\coNP\). Short proof strings give polynomial certificates for
membership in \(\TAUT\). Conversely, when \(\TAUT\) is in \(\NP\), its
verifier defines a proof system by outputting the verified formula and
using a fixed tautology on failed inputs. A bound on the length of an
existing proof leaves the time to find it as a separate question
\cite{cookreckhow}.

\subsection{The quotient by polynomial-time predicates}

Let \(S=\{0,1\}^{*}\), let \(R\subseteq\Ftwo^S\) be the ring of
characteristic functions of \(\PT\) languages, and let \(B\) be the
\(\Ftwo\)-span of the characteristic functions of \(\NP\) languages.
Pointwise conjunction stays in \(\NP\), so \(B\) is a ring and contains
\(R\). The quotient \(B/R\) is an \(R\)-module. The class of
\(\chi_L\), for \(L\in\NP\), vanishes exactly when \(L\in\PT\), and the
whole module vanishes exactly when \(\PT=\NP\).

The semantic comparison in Section~\ref{sec:compatibility} is a
surjection with a kernel ideal. The complexity comparison is an
inclusion with a module cokernel. Relating those two objects requires
a specified map and a comparison theorem. The finite disclosure result
applies to an explicitly presented affine system. Any extension to
general NP witness families would need a uniform construction that
preserves membership and controls the cost of producing a witness.

\subsection{Proof and execution}

The local Lean development proves abstract descent, equality reflection,
tracked substitution, and the pointwise disclosure square. Its tracked
substitution statement assumes a coding and the substitution law.
The logical theorems and the resource bounds for
\eqref{eq:proof-translation} enter through the cited sources. The Python
reference computes with exact cochain coordinates. An execution
refinement would certify that a row-reduction result represents the
entire affine fiber, including the inconsistency case.


\printbibliography
\end{document}
